\chapter{Chapter05 — Threshold Function Counting}
%%% generated by the blueprint pipeline from FABL.Chapter05.ThresholdFunctionCounting
\section{Overview}

Declarations of the Lean module \texttt{FABL.Chapter05.ThresholdFunctionCounting}.

\section{Declarations}

\begin{definition}[Polynomial threshold functions]
\lean{FABL.PolynomialThresholdFunction}
\statementsource{Analysis-of-Boolean-Functions-by-Ryan-ODonnell}{pdf-ffbe9c7b3306-p115-b005}
\label{FABL.PolynomialThresholdFunction}
\leanok
\uses{FABL.BooleanFunction, FABL.IsPolynomialThreshold}
For nonnegative integers $n$ and $k$, a polynomial threshold function on $n$ coordinates
of order $k$ is a Boolean function on $\\{-1,1\\}^n$ that satisfies the
polynomial-threshold property with parameter $k$.
\end{definition}

\begin{definition}[Linear threshold functions]
\lean{FABL.LinearThresholdFunction}
\statementsource{Analysis-of-Boolean-Functions-by-Ryan-ODonnell}{
  pdf-ffbe9c7b3306-p044-b005,
  pdf-ffbe9c7b3306-p113-b003,
  pdf-ffbe9c7b3306-p241-b003
}
\label{FABL.LinearThresholdFunction}
\leanok
\uses{FABL.BooleanFunction, FABL.IsLinearThreshold}
For a positive integer $n$, a linear threshold function on the $n$-dimensional Boolean
cube is a Boolean function on that cube that is a linear threshold function.
\end{definition}

\begin{definition}[Fourier agreement count]
\lean{FABL.fourierAgreementCount}
\label{FABL.fourierAgreementCount}
\leanok
\uses{FABL.BooleanFunction, FABL.parityFunction}
For a natural number $n$, a Boolean function $f:\{\pm1\}^n\to\{\pm1\}$, and a set
$S\subseteq\{1,\ldots,n\}$, the Fourier agreement count is the number of inputs
$x\in\{\pm1\}^n$ for which $f(x)$ equals the parity function associated with $S$ at $x$.
\end{definition}

\begin{theorem}[Fourier coefficient from agreement count]
\lean{FABL.fourierCoeff_eq_agreementCount}
\label{FABL.fourierCoeff_eq_agreementCount}
\leanok
\uses{FABL.BooleanFunction, FABL.BooleanFunction.toReal, FABL.fourierAgreementCount, FABL.fourierCoeff, FABL.monomial, FABL.parityFunction, FABL.parityFunction_toReal}
\difficulty{5}
Let $n\ge 0$, let $f$ map the $n$-dimensional sign cube to $\{-1,1\}$, and let $S$ be a
set of coordinates. If $A$ is the number of inputs on which $f$ agrees with the parity
function on $S$, then the Fourier coefficient of the real-valued version of $f$ at $S$
is $\bigl(2A-2^n\bigr)/2^n$.
\end{theorem}

\begin{definition}[Low-degree agreement profile]
\lean{FABL.lowDegreeAgreementProfile}
\statementsource{Analysis-of-Boolean-Functions-by-Ryan-ODonnell}{pdf-ffbe9c7b3306-p114-b004}
\label{FABL.lowDegreeAgreementProfile}
\leanok
\uses{FABL.BooleanFunction, FABL.fourierAgreementCount, FABL.lowDegreeFourierFamily}
Let $n,k$ be nonnegative integers and let $f$ be a Boolean function on the
$n$-dimensional Boolean cube. The low-degree agreement profile assigns to each subset
$S\subseteq\{1,\ldots,n\}$ with $|S|\le k$ the number of inputs on which $f$ agrees with
the Fourier character indexed by $S$, viewed as an integer from $0$ through $2^n$.
\end{definition}

\begin{theorem}[Injectivity of the polynomial threshold profile]
\lean{FABL.lowDegreeAgreementProfile_injective_on_polynomialThreshold}
\proofsource{Analysis-of-Boolean-Functions-by-Ryan-ODonnell}{pdf-ffbe9c7b3306-p116-b003}
\label{FABL.lowDegreeAgreementProfile_injective_on_polynomialThreshold}
\leanok
\uses{FABL.PolynomialThresholdFunction, FABL.eq_of_isPolynomialThreshold_of_fourierCoeff_eq, FABL.fourierAgreementCount, FABL.fourierCoeff_eq_agreementCount, FABL.lowDegreeAgreementProfile, FABL.lowDegreeFourierFamily, FABL.mem_lowDegreeFourierFamily}
\difficulty{4}
For all natural numbers $n$ and $k$, the map sending a polynomial threshold function on
the $n$-dimensional Boolean cube whose representing polynomial has Fourier degree at
most $k$ to its low-degree agreement profile for degrees at most $k$ is injective. Thus,
two such functions with the same agreement count for every parity function indexed by a
subset of size at most $k$ must be equal.
\end{theorem}

\begin{theorem}[Injectivity of the agreement profile]
\lean{FABL.lowDegreeAgreementProfile_injective_on_linearThreshold}
\proofsource{Analysis-of-Boolean-Functions-by-Ryan-ODonnell}{
  pdf-ffbe9c7b3306-p114-b002,
  pdf-ffbe9c7b3306-p114-b003
}
\label{FABL.lowDegreeAgreementProfile_injective_on_linearThreshold}
\leanok
\uses{FABL.LinearThresholdFunction, FABL.eq_of_isLinearThreshold_of_fourierCoeff_eq, FABL.fourierAgreementCount, FABL.fourierCoeff_eq_agreementCount, FABL.lowDegreeAgreementProfile, FABL.lowDegreeFourierFamily, FABL.mem_lowDegreeFourierFamily}
\difficulty{4}
For every natural number $n$, the map sending a linear threshold function on the
$n$-dimensional Boolean cube to its low-degree agreement profile for degrees at most $1$
is injective. In other words, two such functions with the same number of agreements with
each parity function indexed by a subset of size at most $1$ must be equal.
\end{theorem}

\begin{theorem}[Polynomial threshold function count bound]
\lean{FABL.natCard_polynomialThresholdFunction_le_profileCount}
\label{FABL.natCard_polynomialThresholdFunction_le_profileCount}
\leanok
\uses{FABL.PolynomialThresholdFunction, FABL.lowDegreeAgreementProfile, FABL.lowDegreeAgreementProfile_injective_on_polynomialThreshold, FABL.lowDegreeFourierFamily}
\difficulty{3}
For natural numbers $n$ and $k$, the number of Boolean functions on $n$ coordinates that
admit a polynomial threshold representation of Fourier degree at most $k$ is at most
$(2^n+1)^{|\mathcal{F}_{n,k}|}$, where $\mathcal{F}_{n,k}$ is the family of subsets of
$\{1,\ldots,n\}$ of size at most $k$.
\end{theorem}

\begin{theorem}[Counting linear threshold functions]
\lean{FABL.natCard_linearThresholdFunction_le_profileCount}
\label{FABL.natCard_linearThresholdFunction_le_profileCount}
\leanok
\uses{FABL.LinearThresholdFunction, FABL.lowDegreeAgreementProfile, FABL.lowDegreeAgreementProfile_injective_on_linearThreshold, FABL.lowDegreeFourierFamily}
\difficulty{3}
For every $n\in\mathbb N$, the number of linear threshold functions on $\{−1,1\}^n$ is
at most $(2^n+1)^{n+1}$.
\end{theorem}

\begin{theorem}[Successive powers of two]
\lean{FABL.two_pow_add_one_le_next_two_pow}
\label{FABL.two_pow_add_one_le_next_two_pow}
\leanok
\difficulty{3}
For every natural number $n$, $2^n+1\le 2^{n+1}$.
\end{theorem}

\begin{theorem}[Nat card linear threshold function le two pow sq]
\lean{FABL.natCard_linearThresholdFunction_le_two_pow_sq}
\statementsource{Analysis-of-Boolean-Functions-by-Ryan-ODonnell}{pdf-ffbe9c7b3306-p134-b006}
\label{FABL.natCard_linearThresholdFunction_le_two_pow_sq}
\leanok
\uses{FABL.LinearThresholdFunction, FABL.card_lowDegreeFourierFamily_eq_sum_choose, FABL.lowDegreeFourierFamily, FABL.natCard_linearThresholdFunction_le_profileCount, FABL.two_pow_add_one_le_next_two_pow}
\difficulty{4}
% dataset-pipeline: stub -- prose pending build_informal_statements.py
Statement pending regeneration.
\end{theorem}

\begin{theorem}[Leading term bound for low degree subsets]
\lean{FABL.card_lowDegreeFourierFamily_le_leadingTerm}
\label{FABL.card_lowDegreeFourierFamily_le_leadingTerm}
\leanok
\uses{FABL.card_lowDegreeFourierFamily_eq_sum_choose, FABL.lowDegreeFourierFamily}
\difficulty{4}
For natural numbers $n$ and $k$ with $n>0$, the number of subsets of an $n$-element set
having size at most $k$ is at most $n^k+k n^{k-1}$.
\end{theorem}

\begin{theorem}[Nat card polynomial threshold function le two pow book bound of pos]
\lean{FABL.natCard_polynomialThresholdFunction_le_two_pow_bookBound_of_pos}
\statementsource{Analysis-of-Boolean-Functions-by-Ryan-ODonnell}{pdf-ffbe9c7b3306-p134-b006}
\label{FABL.natCard_polynomialThresholdFunction_le_two_pow_bookBound_of_pos}
\leanok
\uses{FABL.PolynomialThresholdFunction, FABL.card_lowDegreeFourierFamily_le_leadingTerm, FABL.lowDegreeFourierFamily, FABL.natCard_polynomialThresholdFunction_le_profileCount, FABL.two_pow_add_one_le_next_two_pow}
\difficulty{5}
% dataset-pipeline: stub -- prose pending build_informal_statements.py
Statement pending regeneration.
\end{theorem}

\begin{theorem}[Polynomial threshold function book bound]
\lean{FABL.natCard_polynomialThresholdFunction_le_two_pow_bookBound}
\label{FABL.natCard_polynomialThresholdFunction_le_two_pow_bookBound}
\leanok
\uses{FABL.BooleanFunction, FABL.PolynomialThresholdFunction, FABL.Sign, FABL.natCard_polynomialThresholdFunction_le_two_pow_bookBound_of_pos}
\difficulty{4}
For natural numbers $n$ and $k$, the number of Boolean functions on $n$ coordinates that
admit a polynomial threshold representation of Fourier degree at most $k$ is at most
$2^{n^{k+1}+(2k+1)n^k+1}$.
\end{theorem}
